\subsection{关于紧支集测度的积分}

设 \(\mu\) 是 X 上的一个测度，其支集 \(S = \operatorname{Supp}(\mu)\) 是紧的；开集 \(X - S\) 是可忽略的（§2, No.~2, 命题 5）。对每个函数 \(\mathbf{f}\)（取值于向量空间 \(\mathbf{F}\) 或 \(\overline{\mathbf{R}}\) 中），函数 \(\mathbf{f}\) 与 \(\mathbf{f}\varphi_{\mathrm{S}}\) 因而是等价的（§2, No.~4）；为使 \(\mathbf{f}\) 为 \(\mu\)-可积的（当 \(\mathbf{F}\) 是 Banach 空间时），其充要条件是 \(\mathbf{f}\varphi_{\mathrm{S}}\) 可积，此时有（No.~1）

\[
\int \mathbf {f} d \mu = \int \mathbf {f} \varphi_ {\mathrm{S}} d \mu .\tag{21}
\]

此外，若 \(\mathbf{f}\) 在 \(S\) 上有界，则由 (20) 得出

\[
\left| \int \mathbf {f} d \mu \right| \leqslant \| \mu \| \cdot \sup _ {x \in S} | \mathbf {f} (x) |.\tag{22}
\]

特别地，若 \(\mathbf{f}\) 在 \(X\) 上连续，则 \(\mathbf{f}\) 是 \(\mu\)-可积的，因为 \(\mathbf{f}h \in \mathcal{K}(X;F)\)（其中 \(h \in \mathcal{K}(X;R)\) 是在 \(S\) 上等于 1 的任一函数）（Ch. III, §1, No.~2, 引理 1）。更确切地说：

命题 14. --- 设 X 是一个局部紧空间，F 是一个不退化为 0 的 Banach 空间；给 X 到 F 中的全体连续映射所成的空间 \(\mathcal{C}(\mathrm{X};\mathrm{F})\) 赋以紧收敛拓扑。为使 X 上的测度 \(\mu\) 满足：线性映射 \(\mathbf{f} \mapsto \int \mathbf{f} d\mu\)（\(\mathcal{K}(\mathrm{X};\mathrm{F})\) 到 F 中的）可延拓为 \(\mathcal{C}(\mathrm{X};\mathrm{F})\) 到 F 中的连续线性映射，其充要条件是 \(\operatorname{Supp}(\mu)\) 为紧的；这样的延拓是唯一的，且与 No.~1 中定义的积分一致。

我们刚刚看到，若 \(\mu\) 具有紧支集，则积分 \(\int \mathbf{f}d\mu\) 对每个函数 \(\mathbf{f}\in \mathcal{C}(\mathrm{X};\mathrm{F})\) 都有定义，且映射 \(\mathbf{f}\mapsto \int \mathbf{f}d\mu\)（\(\mathcal{C}(\mathrm{X};\mathrm{F})\) 到 \(\mathrm{F}\) 中的）关于紧收敛拓扑是连续的。反之，假设 \(\mathbf{f}\mapsto \int \mathbf{f}d\mu\) 在 \(\mathcal{K}(\mathrm{X};\mathrm{F})\) 中关于紧收敛拓扑是连续的。那么，存在一个紧集 \(K\subset X\) 和一个数 \(a > 0\)，使得 \(|\mu (\mathbf{f})|\leqslant a\cdot \sup_{x\in K}|\mathbf{f}(x)|\) 对每个函数 \(\mathbf{f}\in \mathcal{K}(\mathrm{X};\mathrm{F})\) 成立；特别地，若 \(\mathbf{g}\in \mathcal{K}(\mathrm{X};\mathrm{F})\) 的支集与 \(K\) 不相交，则 \(\mu (\mathbf{g}) = 0\)。取 \(\mathbf{g} = h\mathbf{a}\)，其中 \(\mathbf{a}\neq 0\) 是 \(\mathrm{F}\) 中的一个向量，而 \(h\in \mathcal{K}(\mathrm{X};\mathbf{C})\)，我们看出，\(\mu (h) = 0\) 对每个函数 \(h\in \mathcal{K}(\mathrm{X};\mathbf{C})\)（其支集与 \(K\) 不相交者）都成立，这就证明了 \(\operatorname {Supp}(\mu)\subset K\)。最后，延拓的唯一性由以下事实得出：\(\mathcal{K}(\mathrm{X};\mathrm{F})\) 在 \(\mathcal{C}(\mathrm{X};\mathrm{F})\) 中关于紧收敛拓扑稠密（Ch. III, §1, No.~2, 命题 4）。

命题 14 使我们可以把 X 上具有紧支集的测度与它在 \(\mathcal{C}(\mathrm{X};\mathbf{C})\) 上的连续延拓等同起来。于是，X 上具有紧支集的测度所成之集可以与对偶空间 \(\mathcal{C}'(\mathrm{X};\mathbf{C})\)（即 Hausdorff 局部凸空间 \(\mathcal{C}(\mathrm{X};\mathbf{C})\) 的对偶）等同。回忆起 \(\mathcal{C}(\mathrm{X};\mathbf{C})\) 是完备的（GT, X, §1, No.~6, 定理 2 的推论 3），但它未必是桶式的（习题 17）。然而，若 X 在无穷远处可数，因而是递增紧集序列 \(\mathrm{K}_n\)（满足 \(\mathrm{K}_n \subset \overset{\circ}{\mathrm{K}}_{n+1}\)）之并，则 \(\mathcal{C}(\mathrm{X};\mathbf{C})\) 的拓扑可由可数的半范数族 \(p_n(f) = \sup_{x \in \mathrm{K}_n} |f(x)|\) 定义，因此在这一情形 \(\mathcal{C}(\mathrm{X};\mathbf{C})\) 是一个 Fréchet 空间。由此，设 \(\mathfrak{S}\) 是 \(\mathcal{C}(\mathrm{X};\mathbf{C})\) 的由有界集组成的任一覆盖，则空间 \(\mathcal{C}'(\mathrm{X};\mathbf{C})\) 关于 \(\mathfrak{S}\)-拓扑是拟完备的（TVS, III, §4, No.~2, 定理 1 的推论 4）。

我们在 \(\mathcal{C}^{\prime}(\mathrm{X};\mathbf{C})\) 上主要考虑紧收敛拓扑（即在 \(\mathcal{C}(\mathrm{X};\mathbf{C})\) 的紧子集上一致收敛的拓扑）。回忆起 \(\mathcal{C}(\mathrm{X};\mathbf{C})\) 的相对紧子集 H 由下列性质刻画（GT, X, §2, No.~5, 定理 2 的推论 3）：

\(1^{\circ}\) H 等度连续；

\(2^{\circ}\) 对每个 \(x \in \mathrm{X}\)，集 \(\mathrm{H}(x)\)——诸 \(f(x)\)（其中 \(f\) 取遍 \(\mathrm{H}\)）所成之集——在 \(\mathbf{C}\) 中有界。

命题 15. --- 设 X 是一个局部紧空间，且对每个 \(x \in X\)，设 \(\varepsilon_{x}\) 是点 x 处的 Dirac 测度。映射 \(x \mapsto \varepsilon_{x}\)（从 X 到 \(\mathcal{C}'(X; \mathbf{C})\) 中）关于 \(\mathcal{C}'(X; \mathbf{C})\) 上的紧收敛拓扑是连续的。

考虑 \(\varepsilon_{x_0}\) 在 \(\mathcal{C}'(\mathrm{X};\mathbf{C})\) 中关于这拓扑的一个邻域，我们可以假定它是这样定义的：取一个数 \(\delta > 0\)、一个紧子集 \(\mathrm{H}\)（属于 \(\mathcal{C}(\mathrm{X};\mathbf{C})\)），并考虑测度 \(\mu\)（\(\mathrm{X}\) 上具有紧支集者）所成的集，其中 \(|\mu(f) - \varepsilon_{x_0}(f)| \leqslant \delta\) 对每个函数 \(f \in \mathrm{H}\) 成立。由于 \(\mathrm{H}\) 等度连续，存在 \(\mathrm{U}\)（\(x_0\) 在 \(\mathrm{X}\) 中的一个邻域），使得关系 \(f \in \mathrm{H}\) 蕴涵 \(|f(x) - f(x_0)| \leqslant \delta\) 对所有 \(x \in \mathrm{U}\) 成立，它也可写成 \(|\varepsilon_x(f) - \varepsilon_{x_0}(f)| \leqslant \delta\)，这就证明了命题。(*)

命题 16. --- 设 K 是 X 的一个紧子集，L 是 X 上支集含于 K 的测度 \(\mu\) 所成的向量空间。在 L 上，由紧收敛拓扑 \(\mathcal{T}\)（\(\mathcal{C}'(X; \mathbf{C})\) 上的）所诱导的拓扑与严格紧收敛拓扑 \(\mathcal{T}'\)（\(\mathcal{M}(X; \mathbf{C})\) 上的）（Ch. III, §1, No.~10）一致。

显然，在 L 上，由 T 诱导的拓扑细于由 \(T'\) 诱导的拓扑。反之，设 H 是 \(\mathcal{C}(\mathrm{X};\mathbf{C})\) 的一个紧子集，h 是 \(\mathcal{K}(\mathrm{X};\mathbf{C})\) 中在 K 上等于 1 的一个函数。显然，由函数 fh（其中 f 取遍 H）所成的集 \(H'\) 在 \(\mathcal{K}(\mathrm{X};\mathbf{C})\) 中是严格紧的，且对每个测度 \(\mu\in L\)，\(\mu(f)=\mu(fh)\) 对每个函数 \(f\in H\) 成立，由此即得结论。

推论 1. --- 对 X 的每个紧子集 K 与每个数 a \textgreater{} 0，X 上的测度 \(\mu\) 所成的集 B，其中 \(\operatorname{Supp}(\mu) \subset K\) 且 \(\|\mu\| \leqslant a\)，它是 \(\mathcal{C}'(X; \mathbf{C})\) 的一个等度连续子集，并且它关于紧收敛拓扑 T 是紧的。

因为，设 H 是 \(\mathcal{C}(\mathrm{X};\mathbf{C})\) 的一个由在 K 上一致有界的函数组成的子集；存在一个数 c\textgreater0，使得 \(|\mu(f)|\leqslant c\cdot\|\mu\|\leqslant ac\) 对每个函数 \(f\in H\) 与每个测度 \(\mu\in B\) 成立（由 (22)）；因此 \(B\subset acH^{\circ}\) 在对偶空间 \(\mathcal{C}'(\mathrm{X};\mathbf{C})\)（\(\mathcal{C}(\mathrm{X};\mathbf{C})\) 的对偶）中成立，这就证明了 B 的等度连续性；B 关于 T 的紧性由以下事实得出：在 B 上，T 与淡拓扑诱导出同一拓扑（命题 16 与 Ch. III, §1, No.~10, 命题 17），并且 B 是淡紧的（Ch. III, §1, No.~9, 命题 15 的推论 2 与 §2, No.~2, 命题 6）。

推论 2. --- 每个具有紧支集的测度（相应地，每个具有紧支集的正测度）\(\mu\) 都属于 \(\mathcal{C}'(\mathrm{X};\mathbf{C})\) 中（关于紧收敛拓扑 \(\mathcal{T}\)）由测度（相应地

正测度）——其支集有限、含于 \(\operatorname{Supp}(\mu)\) 且其范数等于 \(\| \mu \|\)——所成之集的闭包。

因为，在测度 \(\nu\)（满足 \(\operatorname{Supp}(\nu) \subset \operatorname{Supp}(\mu)\) 且 \(\|\nu\| \leqslant \|\mu\|\)）所成的集 B 上，由淡拓扑诱导的拓扑与由 T 诱导的拓扑相同，因此该推论由 Ch. III, §2, No.~4, 定理 1 的推论 2 与 3 得出。

